\title{Physics-Constrained Reconstruction of Saturated \texorpdfstring{$dq$}{dq} Flux-Linkage Maps from Magnetic Co-Energy}
\author{}
\date{7 October 2026}

\newcommand{\PaperAbstract}{%
Stationary electrical measurements provide flux-linkage samples, not measured
magnetic energy. This learning paper develops magnetic co-energy as a latent
potential for a conservative storage subsystem in conjugate current coordinates.
An even potential enforces flux reflection and reciprocal differential
inductance, while its varying Hessian represents saturation. Direct regularized
gradient fitting avoids integrating noisy samples along arbitrary paths.
Torque adds a different observation: at nonzero stator current it measures
exactly the flux projection normal to current. A local minimum-change
correction selects one member of an underdetermined solution line.
Smooth radial potentials remain invisible even with reciprocity, reflection,
and positive added curvature; fixing the energy gauge does not remove them.
Electrical and torque observation rows can complement each other, but ideal
rank does not resolve unqualified calibration or loss parameters.
Reproducible synthetic tests separate flux accuracy, derivative quality,
admissibility, torque observability, and noise conditioning.
A real-data stress test retains the disagreement between symmetry and
circulation equivalents despite exact P1 integration. A separate multi-speed
CAN comparison exhibits a strong directional residual under an explicit
unit/sign hypothesis. Recorder metadata do not qualify CAN torque as
electromagnetic or shaft torque, so that comparison demonstrates conditional
consistency rather than experimental validation or correction accuracy.
No real-data magnetic potential or causal error correction is identified.
Structured residuals remain evidence to preserve, not errors to delete.}

% Disable either component if the paper does not need it.
\newif\ifshowglossary
\showglossarytrue
\newif\ifshowbibliography
\showbibliographytrue
